% PROPOSAL: the author rewrites this section. It states the claim as the
% generated numbers support it. Framing and emphasis are the author's decision.
\section{Introduction}

The Henry Hub--TTF spread is the price of moving a unit of gas from the United
States to Europe, and liquefied natural gas (LNG) carriers are the means by which it
moves. A laden carrier that leaves Sabine Pass arrives at Rotterdam roughly fourteen
to eighteen days later. Between departure and arrival its cargo is committed, is
visible to anyone with an AIS receiver, and is absent from every published supply
statistic. It is therefore natural to expect vessel positions to lead the spread.
This expectation is widely held and is the basis of a commercial data industry.

This paper tests it and finds no support. An efficient-markets view predicts that
outcome, so the null itself is not the contribution. The contribution lies in how
the null is established, and in two construction problems encountered on the way.

Three features of the design distinguish the result from a mere absence of findings.
First, every target, control set, lag, sign and acceptance bar was written into a
dated, append-only decision log before the corresponding fit, so no choice reported
below was made after seeing its outcome. Where a pre-registered sign was falsified,
it is reported as falsified. Second, the inference respects the structure of the
data. Weekly changes over a four-week horizon overlap, so all standard errors are
Newey--West with the lag fixed in advance at \hacLagRule{} weeks. Training rows whose
target reaches into the test week are purged. The three follow-up mechanism tests
are held to a Bonferroni bar and must replicate their sign on a holdout that was not
examined until they had passed on the discovery sample. Third, the study bounds
what it could have found. With \powerN{} weekly observations the design detects
effects of at least \powerFloor{} of residual variance, and the largest observed
anywhere is \bMaxRsqCov. The result excludes effects large enough to trade and is
silent on smaller ones.

Two further findings are methodological and apply to any feature built from
vessel-tracking data. The natural unit of observation, a voyage, does not exist
until both of its endpoints have been observed, so a statistic computed over
completed voyages embeds the future unless the population is explicitly bounded. At
a 2020 as-of date a straightforward loader returned \leakNaive{} legs, of which
\leakFutureClosed{} (\leakFutureShare) had been closed by arrivals that had not yet
occurred. Separately, the availability of a feature can change across data vintages.
The same loader attached \leakDeclared{} destination declarations to 2020 voyages
from a data feed that began in 2026. A date filter on observations cannot catch
this, because the individual rows are correctly dated; the table itself should not
exist at that vintage.

The paper proceeds as follows. Section~\ref{sec:data} describes the pipeline and
the decade-scale reconstruction. Section~\ref{sec:panel} explains how the signal
panel is built point-in-time. Section~\ref{sec:prereg} describes the
pre-registration protocol and its limits. Section~\ref{sec:parta} reports the
physical nowcasts. Section~\ref{sec:partb} reports the spread tests.
Section~\ref{sec:power} bounds their power. Sections~\ref{sec:discussion}
and~\ref{sec:limitations} interpret and qualify the results.
